\section{From relation products to the state lower bound}
\label{sec:separation}

We now connect the two representations. The use of the full relation
alphabet makes this step elementary: every relation occurs as a single
letter, and two additional letters select any one entry of a product.
This is the context-separation property of one-way
liveness~\cite{AdeogunKapoutsis2026}.

\begin{proposition}
\label{prop:small-nfa}
For $h\ge2$, the language $\OWL_h$ has a nondeterministic automaton with
$h+3$ states, with no left moves.
\end{proposition}

\begin{proof}
Use the $h$ elements of $H=[h]$ as path states, together with three distinct
states $q_{\mathrm{in}},q_{\mathrm{acc}},q_{\mathrm{rej}}$. Start in
$q_{\mathrm{in}}$ on the left endmarker and move right into any path state
$p\in H$. At a letter $R\in\R{H}$ in path state $p$, move right into any
$q$ with $(p,q)\in R$. If there is no such $q$, enter
$q_{\mathrm{rej}}$ by a stay move. At the right endmarker in any path
state, enter $q_{\mathrm{acc}}$ by a stay move. Declare only
$q_{\mathrm{acc}}$ accepting. The rejecting state stays forever, and all
unused rules may enter it by a stay move.

An accepting computation on $R_1\cdots R_\ell$ specifies states
$p_0,\ldots,p_\ell$ with $(p_{j-1},p_j)\in R_j$ for every $j$.
Such states exist exactly when $R_1\cdots R_\ell\ne0$.
For the empty word, the first right move goes directly from the left
marker to the right marker in a path state, and the next step accepts.
This agrees with $\Id{H}\ne0$. All transitions respect the boundaries,
all $h+3$ states are counted, and acceptance always follows a positive
transition.
\end{proof}

\begin{lemma}
\label{lem:syntactic-quotient}
Let $H$ be a nonempty finite set. Suppose a deterministic automaton
recognizing nonzero relation products over $\R{H}$ has a symbol-diagram
representation in $\D{k}$ as in Theorem~\ref{thm:machine-diagrams}.
Let $S\subseteq\D{k}$ be the submonoid generated by the ordinary-symbol
diagrams, including the empty product. Then there is a unital surjection
$S\to\R{H}$ sending each ordinary-symbol diagram to its relation letter.
\end{lemma}

\begin{proof}
For a body word $w$, let $d(w)\in S$ be its diagram product and let
$r(w)\in\R{H}$ be its relation product. Suppose $d(u)=d(v)$. For any
body words $x,y$, associativity gives
\[
 d_{\mathrm L}\,d(x)d(u)d(y)\,d_{\mathrm R}
 =d_{\mathrm L}\,d(x)d(v)d(y)\,d_{\mathrm R},
\]
where $d_{\mathrm L},d_{\mathrm R}$ are the fixed marker diagrams.
The recognition rule therefore gives the same membership decision for
$xuy$ and $xvy$.

If $r(u)\ne r(v)$, choose $(i,j)$ belonging to exactly one of them.
Take $x$ to be the single letter $\Id{\{i\}}$ and $y$ the single letter
$\Id{\{j\}}$. For any relation $A$,
\[
 \Id{\{i\}}A\Id{\{j\}}
 =\begin{cases}
   \{(i,j)\}, & (i,j)\in A,\\
   0, & (i,j)\notin A.
  \end{cases}
\]
Exactly one of the two surrounded words would belong to the language,
a contradiction. Thus $d(u)=d(v)$ implies $r(u)=r(v)$, and
$\phi(d(w))=r(w)$ is well-defined on $S$.

Concatenation proves that $\phi$ preserves multiplication. The empty
word gives $\phi(\id)=\Id{H}$, including when a nonempty word also
has diagram product $\id$. Every relation is the image of its one-letter
diagram, so $\phi$ is surjective.
\end{proof}

\begin{proof}[Proof of Theorem~\ref{thm:main}]
The source-state assertion is Proposition~\ref{prop:small-nfa}.
Let an arbitrary $s$-state deterministic two-way automaton recognize
$\OWL_h$ on all finite words. Under initial-configuration acceptance,
Theorem~\ref{thm:machine-diagrams} gives a representation in
$\D{k}$ with $k=4(s+1)^2$. Lemma~\ref{lem:syntactic-quotient} supplies
a unital surjection from the generated submonoid onto $\R{[h]}$.
Corollary~\ref{cor:divisor-degree} then gives
\[
 4(s+1)^2=k\ge\Lfun(h)=2^{\lfloor(h-2)/31\rfloor}.
\]
With positive-transition acceptance, use the additional initial copy in
Theorem~\ref{thm:machine-diagrams}; its degree is $4(s+2)^2$.
This proves both stated bounds.
\end{proof}

\begin{remark}
The auxiliary vertices and wires used in the matching representation
are an algebraic encoding of a computation graph, not additional memory
granted to the deterministic simulator. The lower bound is proved for
every machine in the ordinary model stated in the introduction. The
family uses a growing finite alphabet, as permitted by the required
alphabet-independent constants; the proof makes no separate
fixed-alphabet claim.
\end{remark}
